%% Created by Maple V Release 4 (IBM INTEL NT)
%% Source Worksheet: trucks-6.txt
%% Generated: Fri Jul 04 13:11:54 1997
\documentclass{article}
\usepackage{maple2e}
\DefineParaStyle{Heading 1}
\DefineParaStyle{Heading 2}
\DefineParaStyle{Heading 3}
\DefineParaStyle{Maple Output}
\DefineParaStyle{Title}
\DefineCharStyle{2D Math}
\DefineCharStyle{2D Output}
\begin{document}
\begin{maplegroup}
\begin{Title}
Social Structures and Speeding Trucks: A Simulation with Experimentation
\vskip .25in
Michael~A. Faia
\vskip .25in
\end{Title}

This study develops a Maple simulation and exploits opportunities for
experimenting with it; additional examples are available at

http://morton.wm.edu/public\_html/index.html.

The simulation involves truck transportation, incentives for speeding,
and the efficacy of various ways of controlling speeding. After
setting several initial structural conditions, we show how these
conditions make possible an estimate of operating costs per trip as a
function of speed. Next, we suggest ways in which the total cost
function, operating as a feedback mechanism, might influence a
driver's incentives for speeding or slowing down, or a firm's
incentives for encouraging or discouraging speeding. Finally, a series
of experiments, only partially shown herein and involving the
probability and severity of traffic citations, suggests why maximum
deterrence of speeding requires a judicious combination of sanction
severity and sanction probability.

Reading Ouellet (1994), we learn that the typical trucker logs 60 to
70 hours per week; that he works ``at a steady clip" whether paid by
the hour or by the load; that there are subtle pressures that lead to
excessive speed, e.g., ``drivers often jockey one another for position
on the road and by the loading dock, as a minute or two lead to the
next job ... can save an hour or two of waiting time." We therefore
seek conditions under which a speed in excess of 65 to 70 miles per
hour would appear to be ``rational." Are there circumstances under
which, say, 80 m.p.h. would seem rational? How would law enforcement
attempt to manipulate these circumstances? Here, of course, the word
rational means cost-minimizing, or profit-maximizing.

The user, early in this simulation, may experiment with the first four
input values, k1, x[1], k3, and k4. Later on, we encounter additional
opportunities for experimentation based on further theoretical
considerations.

\begin{Heading 3}
Structure
\end{Heading 3}

First, we insert a baseline operating cost, in dollars per mile, for a
given type of truck. By baseline, we mean all costs involved in paying
for and maintaining a truck, worked out on a per-mile basis, but
without considering ways in which costs are raised by driver behavior,
e.g., high operating speeds, mechanical abuse, etc., or by traffic
citations.

Standard cost per mile, in dollars:

\begin{mapleinput}
\mapleinline{active}{1d}{k1 := 1/3;}{%
}
\end{mapleinput}

\mapleresult
\begin{maplelatex}
\[
{\it k1} := {\displaystyle \frac {1}{3}}
\]
\end{maplelatex}

\end{maplegroup}
\begin{maplegroup}
Truck operating costs increase exponentially with average speed, due
to increased fuel consumption and increased depreciation. We estimate
this additional per/mile cost to be a certain multiplier of a dollar:
the speed in m.p.h., called v, which we will first adjust with an
exponent and then divide by, say, 1000 or some other large
``standard." This latter value itself may need to be adjusted from
time to time to reflect changing costs of fuel, depreciation, etc. We
begin with an equation such as that given immediately below as a
plausible way of capturing this variable.

\begin{mapleinput}
\mapleinline{active}{1d}{x[1] := (v^(8/5))/(4 * (10^3)); # additional
velocity factor}{%
}
\end{mapleinput}

\mapleresult
\begin{maplelatex}
\[
{x_{1}} := {\displaystyle \frac {1}{4000}} \,v^{8/5}
\]
\end{maplelatex}

\end{maplegroup}
\begin{maplegroup}
The next several expressions should be fairly clear. They, too, are
subject to change. We begin by considering the situation of drivers
who have an hourly pay rate; such drivers, according to Ouellet, are
in a very different situation from those who receive piece-work
payment. The hourly wage, initialized below, assumes that for at least
part of a lengthy trip, two drivers must be paid, so that the value
assigned to k3 is a weighted average.

\begin{mapleinput}
\mapleinline{active}{1d}{k3 := 3715/100; # drivers' hourly wage}{%
}
\end{mapleinput}

\mapleresult
\begin{maplelatex}
\[
{\it k3} := {\displaystyle \frac {743}{20}}
\]
\end{maplelatex}

\end{maplegroup}
\begin{maplegroup}
\begin{mapleinput}
\mapleinline{active}{1d}{k4 := 600; # trip length, miles}{%
}
\end{mapleinput}

\mapleresult
\begin{maplelatex}
\[
{\it k4} := 600
\]
\end{maplelatex}

\end{maplegroup}
\begin{maplegroup}
Now we obtain a provisional total cost, for trips of this type. The
total cost has two components, as in a Cobb-Douglas function (Dowling
1990:213-14): non-labor cost and labor cost. The first component
equals the length of the trip multiplied by the sum of k1 and x[1]:

\begin{mapleinput}
\mapleinline{active}{1d}{cn := k4*(k1 + x[1]); # non-labor cost,
entire trip}{%
}
\end{mapleinput}

\mapleresult
\begin{maplelatex}
\[
{\it cn} := 200 + {\displaystyle \frac {3}{20}} \,v^{8/5}
\]
\end{maplelatex}

\end{maplegroup}
\begin{maplegroup}
\begin{mapleinput}
\mapleinline{active}{1d}{plot(cn, v=50..90, thickness = 3, labels =
[speed, cost], title = `Fig. 1: Non-labor cost as function of
speed`);}{%
}
\end{mapleinput}

\end{maplegroup}
\begin{maplegroup}
The second cost component equals the driver's hourly wage multiplied
by the amount of time required for the trip, which is the trip length
divided by the speed:

\begin{mapleinput}
\mapleinline{active}{1d}{cl := k3*(k4/v); # labor cost, entire trip}{%
}
\end{mapleinput}

\mapleresult
\begin{maplelatex}
\[
{\it cl} := {\displaystyle \frac {22290}{v}}
\]
\end{maplelatex}

\end{maplegroup}
\begin{maplegroup}
\begin{mapleinput}
\mapleinline{active}{1d}{plot(cl, v=50..90, thickness = 3, labels =
[speed, cost], title = `Fig. 2: Labor cost as function of speed`);}{%
}
\end{mapleinput}

\end{maplegroup}
\begin{maplegroup}
Examine non-labor costs and labor costs together, as functions of
speed, on a single plot. When we set out to find the speed that
minimizes costs, we shall realize that this is a speed for which
non-labor costs are increasing at the same rate as labor costs are
decreasing.

\begin{mapleinput}
\mapleinline{active}{1d}{plot(\{cn, cl\}, v=50..90, thickness = 3,
labels = [speed, cost], title = `Fig. 3: cn and cl as function of
speed`);}{%
}
\end{mapleinput}

\end{maplegroup}
\begin{maplegroup}
Finally, we obtain the sum of cn and cl, which gives us the total
cost:

\begin{mapleinput}
\mapleinline{active}{1d}{ct := cn + cl; #total cost, entire trip}{%
}
\end{mapleinput}

\mapleresult
\begin{maplelatex}
\[
{\it ct} := 200 + {\displaystyle \frac {3}{20}} \,v^{8/5} +
{\displaystyle \frac {22290}{v}}
\]
\end{maplelatex}

\end{maplegroup}
\begin{maplegroup}
\begin{mapleinput}
\mapleinline{active}{1d}{plot(ct, v=50..90, thickness = 3, labels =
[speed, cost], title = `Fig. 4: Total cost as function of speed`);}{%
}
\end{mapleinput}

\end{maplegroup}
\begin{maplegroup}
\begin{Heading 3}
Dynamics
\end{Heading 3}

In order to analyze the socio-ecological dynamics of the
truck-transportation situation as defined thus far, we begin by
obtaining a derivative as follows:

\begin{mapleinput}
\mapleinline{active}{1d}{Diff(ct, v) = diff(ct, v);}{%
}
\end{mapleinput}

\mapleresult
\begin{maplelatex}
\[
{\frac {\partial }{\partial v}}\,(200 + {\displaystyle \frac {3}{
20}} \,v^{8/5} + {\displaystyle \frac {22290}{v}} )=
{\displaystyle \frac {6}{25}} \,v^{3/5} - {\displaystyle \frac {
22290}{v^{2}}}
\]
\end{maplelatex}

\end{maplegroup}
\begin{maplegroup}
This expression gives the rate of change of cost, relative to speed.
We can use it to find the speed that minimizes total costs; this speed
is often very high---over 80 m.p.h. in this initial simulation---and
it is sometimes considerably higher.

\begin{mapleinput}
\mapleinline{active}{1d}{sols := solve(diff(ct, v) = 0, v): sols[1];
evalf(", 4);}{%
}
\end{mapleinput}

\mapleresult
\begin{maplelatex}
\[
92875^{5/13}
\]
\end{maplelatex}

\begin{maplelatex}
\[
81.42
\]
\end{maplelatex}

\end{maplegroup}
\begin{maplegroup}
Perhaps more realistically, one could ask how much incentive there
would be to increase one's speed from, say, 60 or 70 m.p.h. to any
given higher level. That is, how much would costs for each trip drop
for each mile-per-hour added to one's speed, starting at a relatively
safe speed? Here, we need to solve the derivative for v = 60 and v =
70. Adjacent values may be substituted, for experimental purposes.

\begin{mapleinput}
\mapleinline{active}{1d}{k5 := 60: k6 := 70:}{%
}
\end{mapleinput}

\end{maplegroup}
\begin{maplegroup}
\begin{mapleinput}
\mapleinline{active}{1d}{subs(v = k5, diff(ct, v)): evalf(", 3);
subs(v = k6, diff(ct, v)): evalf(", 3);}{%
}
\end{mapleinput}

\mapleresult
\begin{maplelatex}
\[
-3.38
\]
\end{maplelatex}

\begin{maplelatex}
\[
-1.48
\]
\end{maplelatex}

\end{maplegroup}
\begin{maplegroup}
These solutions imply that there is a strong incentive for going
faster than 60 m.p.h.---\$3.38 can be shaved from total-trip cost for
each additional mile per hour. At higher speeds, incentives are
weaker.

\begin{Heading 3}
Adaptation and Social Control
\end{Heading 3}

In this example, it appears that the tendency to race for loading
docks, as described by Ouellet, could be exploited by unscrupulous
firms as a way of reducing costs: Any factor that causes drivers to
increase speed toward the minimal-cost operating speed, which may be
very high, will reduce costs for the trucking firm. Suppose, however,
that these cost-reduction strategies tend to get drivers and firms
into trouble with the police and highway patrol. This contingency adds
to costs. Estimate these costs on a miles-per-hour basis for all
speeds in excess of what we might call the de facto speed limit---the
speed where the police start writing tickets.

For now, there are two decisions about the sanction system: First,
what is the probability of a speeding ticket on a given trip, when one
exceeds the de facto speed limit? Second, what is the severity of the
sanction, that is, the cost of speeding tickets, for each m.p.h. in
excess of the de facto limit?

Social scientists who write about deterrence usually distinguish among
four aspects of punishment: probability, severity, celerity (speed),
and saliency (awareness of prior warnings or prior instances of
punishment). How could all four of these factors be written into a
simulation? What is the relative impact of, say, a 10 per cent
increase in the probability of punishment versus a plausible increase
in the severity of it? We address these questions in an extended
version of this paper.

\begin{Heading 3}
Sanctions: I
\end{Heading 3}

We begin working toward a resolution of these problems by writing a
new cost function, one that will take into account our costs as
already established but will also allow for increased costs due to
traffic citations. Initially, we assume a fixed probability for
traffic citations, regardless of the degree to which speed exceeds the
de facto speed limit. The assumption is based on the idea that
excessive speeders make special efforts to avoid detection.

\end{maplegroup}
\begin{maplegroup}
\begin{mapleinput}
\mapleinline{active}{1d}{k7 := 1/10: # probability of receiving a
speeding ticket}{%
}
\end{mapleinput}

\end{maplegroup}
\begin{maplegroup}
\begin{mapleinput}
\mapleinline{active}{1d}{k8 := 9: # sanction severity: cost of
speeding ticket, for each m.p.h. excess}{%
}
\end{mapleinput}

\end{maplegroup}
\begin{maplegroup}
\begin{mapleinput}
\mapleinline{active}{1d}{k9 := 75: # the de facto speed limit}{%
}
\end{mapleinput}

\end{maplegroup}
\begin{maplegroup}
\begin{mapleinput}
\mapleinline{active}{1d}{cc := k7*k8*(v - k9); cn; cl;}{%
}
\end{mapleinput}

\mapleresult
\begin{maplelatex}
\[
{\it cc} := {\displaystyle \frac {9}{10}} \,v - {\displaystyle
\frac {135}{2}}
\]
\end{maplelatex}

\begin{maplelatex}
\[
200 + {\displaystyle \frac {3}{20}} \,v^{8/5}
\]
\end{maplelatex}

\begin{maplelatex}
\[
{\displaystyle \frac {22290}{v}}
\]
\end{maplelatex}

\end{maplegroup}
\begin{maplegroup}
\begin{mapleinput}
\mapleinline{active}{1d}{ctt := cn + cl + cc; # total cost, including
impact of sanction system}{%
}
\end{mapleinput}

\mapleresult
\begin{maplelatex}
\[
{\it ctt} := {\displaystyle \frac {265}{2}}  + {\displaystyle
\frac {3}{20}} \,v^{8/5} + {\displaystyle \frac {22290}{v}}  +
{\displaystyle \frac {9}{10}} \,v
\]
\end{maplelatex}

\end{maplegroup}
\begin{maplegroup}
\begin{mapleinput}
\mapleinline{active}{1d}{plot(\{ct, ctt\}, v=k9..90, thickness = 3,
labels = [speed, cost], title = `Fig. 5: Excess operating costs due to
sanction system, I`, titlefont=[TIMES, ROMAN, 9]);}{%
}
\end{mapleinput}

\end{maplegroup}
\begin{maplegroup}
How effective is our sanction system? If the region located between
these two curves, and ranging from the de facto speed limit to, say,
90 m.p.h., were considered a ``braking factor," then we could
integrate through the region as a way of assessing the weight of this
factor.

\begin{mapleinput}
\mapleinline{active}{1d}{Digits := 7: Int(ctt, v = k9..90) = int(ctt,
v); i := int(ctt, v=k9..90):}{%
}
\end{mapleinput}

\mapleresult
\begin{maplelatex}
\[
{\displaystyle \int _{75}^{90}} {\displaystyle \frac {265}{2}}
 + {\displaystyle \frac {3}{20}} \,v^{8/5} + {\displaystyle
\frac {22290}{v}}  + {\displaystyle \frac {9}{10}} \,v\,dv=
{\displaystyle \frac {265}{2}} \,v + {\displaystyle \frac {3}{52}
} \,v^{13/5} + 22290\,{\rm ln}(v) + {\displaystyle \frac {9}{20}
} \,v^{2}
\]
\end{maplelatex}

\end{maplegroup}
\begin{maplegroup}
\begin{mapleinput}
\mapleinline{active}{1d}{evalf(", 6);}{%
}
\end{mapleinput}

\mapleresult
\begin{maplelatex}
\[
9789.82
\]
\end{maplelatex}

\end{maplegroup}
\begin{maplegroup}
\begin{mapleinput}
\mapleinline{active}{1d}{Int(ct, v = k9..90) = int(ct, v); j :=
int(ct, v=k9..90): evalf(", 6);}{%
}
\end{mapleinput}

\mapleresult
\begin{maplelatex}
\[
{\displaystyle \int _{75}^{90}} 200 + {\displaystyle \frac {3}{20
}} \,v^{8/5} + {\displaystyle \frac {22290}{v}} \,dv=200\,v +
{\displaystyle \frac {3}{52}} \,v^{13/5} + 22290\,{\rm ln}(v)
\]
\end{maplelatex}

\begin{maplelatex}
\[
9688.52
\]
\end{maplelatex}

\end{maplegroup}
\begin{maplegroup}
\begin{mapleinput}
\mapleinline{active}{1d}{weight := i - j; evalf(", 5);}{%
}
\end{mapleinput}

\mapleresult
\begin{maplelatex}
\[
{\it weight} := {\displaystyle \frac {405}{4}}
\]
\end{maplelatex}

\begin{maplelatex}
\[
101.25
\]
\end{maplelatex}

\end{maplegroup}
\begin{maplegroup}
\begin{mapleinput}
\end{mapleinput}

\end{maplegroup}
\begin{maplegroup}
This does not seem to be a very heavy weight, but all is relative. How
could we make it heavier? What is the relative impact of sanction
probability and sanction severity on the braking factor? In a later
installment, we shall address this question.

\vskip .15 in
Reference:

\indent\indent Dowling, Edward T. 1990. Calculus for business, economics, and

the social sciences. New York: McGraw-Hill.

\indent\indent Ouellet, Lawrence J. 1994. Pedal to the metal. Philadelphia: Temple

University Press.

\end{maplegroup}
\end{document}
%% End of Maple V Output
